% R3: Dynamic Time Warping (Sanika). Numbers from scripts/run_dtw.py, seed 0; options and
% decision record in DTW.md. Needs: graphicx, booktabs, amsmath.

\section{Dynamic Time Warping}\label{sec:dtw}

The proposal planned to use DTW on load-transition and speed-change segments. The
published recordings contain no such segments: each is a steady 15-minute run at one
speed and one load, cut into fixed-length windows, so there is no variable-length
alignment problem to solve. We therefore use DTW where it can still apply, as the distance
of a $k$-NN classifier between short sequences derived from each window, and ask whether
elastic alignment adds anything.

\textbf{Setup.} Each window is reduced to its amplitude envelope: the RMS of 16-sample
frames on the three gravity-compensated axes, log-scaled, giving a $64\times3$ sequence.
The distance is multivariate DTW with a Sakoe--Chiba band of $\pm$8 frames. Two controls
use the same reference windows: the same envelopes with the band set to 0 (no warping),
and Euclidean $k$-NN on the 126 standardised hand-written features. Evaluation uses the
grouped folds. To keep the run to about a minute on one CPU core, the references are every
12th window of each training recording ($\approx$450) and the queries every 4th window of
each held-out recording (1,584 in total); with 50\% window overlap, neighbouring windows
are near-duplicates.

\begin{table}[t]
  \centering
  \caption{$k$-NN on window envelopes (grouped CV, mean $\pm$ std over 6 folds; recalls
  pooled over folds). Chance macro-F1 is 0.17.}
  \label{tab:dtw}
  \resizebox{\columnwidth}{!}{\begin{tabular}{llrrr}
\toprule
Distance & $k$ & macro-F1 & B1 recall & R recall \\
\midrule
DTW (band $\pm$8) & 1 & 0.28 $\pm$ 0.07 & 0.54 & 0.61 \\
126 hand features & 1 & 0.30 $\pm$ 0.12 & 0.97 & 0.23 \\
no warping (band 0) & 1 & 0.26 $\pm$ 0.09 & 0.55 & 0.54 \\
DTW (band $\pm$8) & 5 & 0.29 $\pm$ 0.09 & 0.53 & 0.59 \\
126 hand features & 5 & 0.30 $\pm$ 0.12 & 0.99 & 0.19 \\
no warping (band 0) & 5 & 0.28 $\pm$ 0.13 & 0.53 & 0.56 \\
\bottomrule
\end{tabular}
}
\end{table}

\textbf{Results} (Table~\ref{tab:dtw}). Warping adds nothing measurable: DTW reaches a
macro-F1 of 0.28 against 0.26 with no warping ($k=1$), a gap well inside the
fold-to-fold spread of 0.07--0.13. Both envelope distances sit just below the hand-written
features (0.30). This matches what the signals are: two windows of a steady-state run
start at arbitrary points of the shaft cycle, so there is no common event for the warping
path to line up, only random fluctuations. The envelope is nonetheless complementary to
the features. It recognises R far better (recall 0.61 against 0.23) without simply
over-predicting it (precision 0.45 against 0.13), and B1 far worse (recall 0.54 against
0.97). The gain is concentrated in the unloaded runs (R recall 0.88 unloaded, 0.33
loaded), although the features also do better there (0.40 against 0.05). We have not
identified what the envelope responds to. It is not the loaded pole-pass modulation: a
6.25\,Hz envelope resolves nothing above 3.1\,Hz, while the loaded spacing is
3.5--10.1\,Hz (Sec.~\ref{sec:separability}).

\textbf{Conclusion.} On steady-state recordings, DTW's alignment brings no benefit, and we
keep it only as a negative-control baseline. Transition segments, the use the proposal had
in mind, would need recordings of speed or load changes, which this dataset does not
include.
